% ===== §11 The Theory and Existing Forms of Everyday Practice =====
\section{The Theory and Existing Forms of Everyday Practice}
\label{p23:sec:practicefield}

\noindent
The paper is about to propose models of everyday aesthetic practice, and it owes, before it proposes anything, a full account of the theory and the practice that already occupy this ground. Everyday life has been theorised, its practices studied, and its aesthetics given a scholarly literature; aesthetic education has its methods; and the East has traditions of everyday practice older and more refined than any Western theory of them. To advance a typology of practice without first setting it within this field would be to reinvent, in ignorance, what has been carefully thought, and to earn the objection that the paper does not know the literature it presumes to extend. This section therefore surveys the field in four registers, the theory of practice, the practical methods of aesthetic education, the sub-field of everyday aesthetics, and the practice traditions of the East, and it closes by marking precisely what these leave undone, which is the space the paper's own models are built to occupy. The survey is not a formality. Each of its registers has reached a limit that the paper's typology is designed to pass, and stating the limits precisely is how the typology earns the right to be proposed.

\subsection{The theory of practice}
\label{p23:sub:field-practicetheory}

The theory of practice supplies the frame within which any account of everyday doing must situate itself. Bourdieu made the habitus its central concept: the durable dispositions, laid down by a social position, that generate a person's practices without being consciously followed, so that practice is the improvisation of an embodied social sense rather than the execution of a rule \parencite{bourdieu1990}. The account has great descriptive power and a limit that its critics fixed precisely. De Certeau, granting Bourdieu's analyses, objected that the habitus finally subordinates practice to the law of reproduction, reducing the lucidly described improvisations of everyday life to the reproduction of the social structures that formed them \parencite{decerteau1984}. Against this he set the distinction his own work is known for. Strategies belong to the powers that produce and order a place, the institutions that circumscribe a proper domain; tactics belong to those who lack such a place and who operate within the order others have produced, poaching on its territory, making do, turning the arrangements of the powerful to unforeseen uses in a creative resistance that is influenced by the imposed order but never wholly determined by it \parencite{decerteau1984}. Everyday life, for de Certeau, is the art of these tactics, the inventive reappropriation by ordinary people of a world produced by others. This gives the paper two of its coordinates and, in the same gesture, two of the limits it must pass. Bourdieu's habitus theorises everyday practice as reproduction, which is, in this paper's terms, the closed cycle, the practice that returns the social order to itself; de Certeau's tactic theorises it as resistance, which is reactive, defined by the imposed order it diverges from. Neither theorises everyday practice as generation, as the endogenous production of new aesthetic value around an unsettled centre. Reproduction and resistance do not exhaust the possibilities of everyday doing, and the possibility they omit, generation, is the one the paper's models are about.

\subsection{The practical methods of aesthetic education}
\label{p23:sub:field-edupractice}

The tradition of aesthetic education, whose theory was set out earlier (\xref{p23:sec:education}), has also its practical methods, and they belong to this survey because the paper's models are models of practice and must be set beside the practices that already claim this ground. That tradition trained perception and sensibility through determinate means: the guided encounter with works of art, the practice of an art oneself, the disciplines of drawing, music, and craft through which Schiller's and Read's cultivation of the integrated person was actually pursued, and, in the schools that took this most seriously, the sensory education of the whole child, as in the Montessori training of the discriminating senses and the Waldorf cultivation of the aesthetic and imaginative life through rhythm, making, and art. The East placed its practical methods differently, in the long apprenticeships of repeated practice, \zh{稽古}, and in the copying of masters through which the arts of calligraphy, painting, and the tea gathering were transmitted, methods in which perception was cultivated not by the appreciation of works but by the disciplined repetition of a making. These methods are real and effective, and the paper does not disparage them. But their unit, as the theory of the tradition already showed, is the individual whose sensibility they refine, and their setting is the reserved one of the lesson, the studio, the practice hall, the guided encounter set apart from ordinary life. What they do not offer is a practice of the everyday as such, conducted in the unreserved flow of ordinary shared living and taking the relation rather than the individual as what it cultivates. The tradition's methods train the person in a space set apart; the paper's models work the relation in the midst of ordinary life.

\subsection{The sub-field of everyday aesthetics}
\label{p23:sub:field-everydayaesthetics}

Everyday aesthetics is now an established sub-field, and the paper must locate itself within it rather than beside it. Its founding move was to establish that the aesthetic extends far beyond art into the textures of ordinary life, and that this domain deserves the attention aesthetics had reserved for the artwork; Saito's studies of the aesthetics of the familiar are its central statement, and the paper has drawn on them throughout \parencite{saito2007,saito2017}. Leddy located the aesthetic value of ordinary life in the qualities by which the everyday becomes, on attention, extraordinary, developing a broadly Deweyan account of the aura of the ordinary \parencite{leddy2012}. Others extended the field toward the practices and politics of everyday life, reading the ordinary through the lens of cultural theory. The sub-field has, in short, established the domain, described its phenomena, and defended its worth, and the present paper stands on that achievement. But it is marked by a limit its own participants have debated, and the limit is the one this paper's central claim is pitched against. The sub-field has largely treated everyday aesthetic experience as the private experience of an individual perceiver, a matter of how a person finds the ordinary beautiful; Saito's account, in its strongest form, holds everyday aesthetic experience to be private and personal in a way that distinguishes it from the public experience of art. What the sub-field has not chiefly done is to theorise everyday aesthetic experience as relationally generated, as the shared production of a perceived world by those who perceive it together, nor to give the mechanism of that generation. It has established that the everyday is aesthetic and left open how the everyday aesthetic is generated between people, which is the question the paper's convergence answered and its models of practice address.

\subsection{The practice traditions of the East}
\label{p23:sub:field-east}

The East theorised everyday aesthetic practice, as practice, earlier and more deeply than the West, and the paper's models owe these traditions an acknowledgement that is more than decorative, since several of the models are their descendants. The way of tea made a discipline of the shared aesthetic occasion, treating the preparation and taking of tea as a practice in which attention, gesture, utensil, and space are composed into a single unrepeatable event, and giving that event the name the paper has already made central, \zh{一期一会}. Dogen, in the \zh{典座教訓}, the instructions for the monastery cook, made the most ordinary practice imaginable, the preparing of meals, into a discipline of awakened attention: the cook is not to regard the ingredients with ordinary eyes or handle them with ordinary feelings, is to bring the same care to the coarsest greens as to the finest cream, and is to work with the nurturing mind that the text names among the three minds of the practice \parencite{dogen1996}. Cooking, in Dogen, is not preparation for practice but practice itself, the everyday activity through which awakening is actualised, and this is the East's clearest statement of the thesis that everyday doing, rightly done, is where the cultivation of perception occurs. Yanagi's philosophy of \zh{民藝}, the ordinary crafts, located beauty in the objects of daily use and in the unselfconscious making of them, a beauty of use that belongs to the ordinary and the anonymous rather than to the signed artwork \parencite{yanagi1972}. These traditions are, in several respects, ahead of the paper, and it learns from them more than it corrects. What it adds to them is not a deeper practice but a theory: the traditions transmit their practices by apprenticeship and example, within frameworks of religious and cultural meaning that the paper does not presume, and they do not, as such, offer the account of why these practices generate, in the terms of a general theory of perception and value, that the paper's convergence supplies. The paper's models aim to state, in a form that does not depend on a particular tradition's framework, the generative structure that these traditions have long enacted.

\subsection{What the field leaves undone}
\label{p23:sub:field-gap}

The survey can now be gathered into the precise statement of what the field leaves for the paper's models to do. The theory of practice has given everyday doing as reproduction and as resistance, but not as generation. The practical methods of aesthetic education have trained the individual's sensibility in a reserved space, but have not offered a practice of the relation in the unreserved flow of ordinary life. The sub-field of everyday aesthetics has established the everyday as aesthetic and treated it largely as the private experience of an individual, but has not theorised its relational generation or given its mechanism. And the practice traditions of the East have enacted the generative cultivation of everyday perception with a depth the West has not matched, but have transmitted it by apprenticeship within particular frameworks rather than stating its structure as such. The space these leave open is exact and it is single: a set of models of everyday aesthetic practice that are generative rather than reproductive or merely resistant, relational rather than individual, conducted in the flow of ordinary shared life rather than in a reserved space, and stated in the general terms of the paper's theory of perception and value rather than within a particular tradition. That space is what the next section is built to occupy, and the models it proposes are answerable, one by one, to the limits this survey has fixed.
